% vi: fileencoding=utf-8 spelllang=pt_br spell
\documentclass[a4paper,12pt,openright,oneside]{book}

\usepackage{microtype}

\usepackage[T1]{fontenc}
\usepackage[utf8]{inputenc}

\usepackage[brazil]{babel}

\usepackage[bookmarks=true]{hyperref}
\hypersetup{
  pdftitle={},
  pdfauthor={},
  pdfkeywords={},
  bookmarksnumbered,
  breaklinks=true,
  urlcolor=blue,
  citecolor=black,
  colorlinks=true,
  linkcolor=black,
}

\usepackage{lmodern}
\usepackage{amsmath}
\usepackage{amssymb}
\usepackage{textcomp}

\usepackage[style=ieee, citestyle=ieee]{biblatex}
\bibliography{bibliografia}

\usepackage[
	per-mode=symbol,
	output-decimal-marker={,},
	separate-uncertainty=true,
]{siunitx}

\usepackage{booktabs}
\usepackage{caption}
\captionsetup[table]{skip=1ex}

\usepackage{tikz}
\usepackage{graphicx}
\graphicspath{{figures/}}

\usepackage{pdfpages}

\usepackage{enumitem}

\usepackage[a4paper,
	top=3cm,
	bottom=2cm,
	left=3cm,
	right=2cm,
	headsep=1cm,
	headheight=14.5pt,
]{geometry}

\usepackage{setspace}

\usepackage{indentfirst}
\setlength{\parindent}{4em}

\usepackage{fancyhdr}
\fancypagestyle{plain}{\fancyhf{}\rhead{\thepage}}

%\usepackage{titlesec}
%\titleformat{\chapter}[display]
%	{\normalfont\sffamily\bfseries\LARGE}
%	{\chaptertitlename\ \thechapter}
%	{3ex}{\Huge}{}
%\titleformat{\section}
%	{\normalfont\sffamily\bfseries\Large}
%	{\thesection}
%	{1em}{}{}
%\titleformat{\subsection}
%	{\normalfont\sffamily\bfseries\large}
%	{\thesubsection}
%	{1em}{}{}

\AtBeginDocument{\addtocontents{toc}{\protect\thispagestyle{empty}}}
\AtBeginDocument{\addtocontents{lof}{\protect\thispagestyle{empty}}}
\AtBeginDocument{\addtocontents{lot}{\protect\thispagestyle{empty}}}

\newcommand{\fronttitle}[1]{\begin{center}{\large\bfseries #1}\end{center}\vskip 2\baselineskip}
\newcommand{\frontitem}[1]{{\small\bfseries #1}}

\usepackage{relsize}
\usepackage[nomain, acronym, nopostdot]{glossaries}
\setacronymstyle{long-sm-short}
\makenoidxglossaries
\newcommand{\newacronymx}[8][]{%
	\newglossaryentry{#2}{
	type=\acronymtype,
	name={{\smaller #3}},
	sort={#3},
	first={\if\relax\detokenize{#4}\relax%
			\emph{#5} ({\smaller #3})%
		\else%
			#4 ({\smaller #3}, \emph{#5})%
		\fi},
	firstplural={\if\relax\detokenize{#7}\relax%
			\emph{#8} ({\smaller #6})%
		\else%
			#7 ({\smaller #6}, \emph{#8})%
		\fi},
	text={{\smaller #3}},
	plural={{\smaller #6}},
	description={\if\relax\detokenize{#4}\relax%
			\emph{#5}%
		\else%
			#4 (\emph{#5})%
		\fi},
		#1}}

% Abreviaturas e acrônimos.
\newacronym{FEEC}{FEEC}{Faculdade de Engenharia Elétrica e de Computação}
\newacronym{UNICAMP}{UNICAMP}{Universidade Estadual de Campinas}

% É possível definir também as formas plurais
\newacronym[plural=CAs, longplural=centros acadêmicos]{CA}{CA}{centro acadêmico}

% Itens não utilizados no texto não aparecerão na Lista de Abreviaturas.
\newacronym{ANUT}{ANUT}{abreviatura não utilizada no texto}

% Abreviaturas em língua diferente da língua principal devem aparecer com a tradução.
\newacronymx{AAL}{AAL}{acrônimo em outra língua}{acronym in another language}{AALs}{acrônimos em outra língua}{acronyms in another language}



\begin{document}

\glsunsetall

\renewcommand{\contentsname}{Sumário}
\renewcommand{\bibname}{Referências}

\pagestyle{empty}

\begin{center}

\raisebox{-5mm}{\makebox[0pt][l]{\begin{tikzpicture}[x=0.08mm, y=0.08mm]
\definecolor{unicampred}{RGB}{218,41,28}
\fill (23.9531, 1.5) arc (3.58332:18.9167:24) -- (68.5957, 26.7897)arc (20.9142:16.5994:265.82) arc (128.381:171.373:10) -- cycle;
\fill (21.5557, 10.5523) arc (26.0833:41.4167:24) -- (55.0392, 52.9179)arc (32.5881:22.0609:144) -- cycle;
\fill (15.8767, 17.998) arc (48.5833:63.9167:24) -- (28.9371, 65.9407)arc (254.14:317.658:10) arc (56.8651:36.5278:55.38) -- cycle;
\fill (7.78063, 22.7038) arc (71.0833:86.4167:24) -- (1.5, 80.0404)arc (89.1258:74.8637:82.58) arc (170.647:236.876:10) -- cycle;
\fill (-1.5, 23.9531) arc (93.5833:108.917:24) -- (-28.6966, 73.1995)arc (115.826:92.2091:68.49) -- cycle;
\fill (-10.5523, 21.5557) arc (116.083:131.417:24) -- (-53.0458, 55.1671)arc (138.056:117.14:75.31) -- cycle;
\fill (-17.998, 15.8767) arc (138.583:153.917:24) -- (-67.7798, 29.6989)arc (162.818:139.707:66.11) -- cycle;
\fill (-22.7038, 7.78063) arc (161.083:176.417:24) -- (-67.7575, 1.5)arc (190.38:172.687:82.03) -- cycle;
\fill (-23.9531, -1.5) arc (183.583:198.917:24) -- (-56.342, -21.714)arc (214.771:200.588:92.45) -- cycle;
\fill (-21.5557, -10.5523) arc (206.083:221.417:24) -- (-38.6627, -36.5413)arc (237.639:225.296:92.84) -- cycle;
\fill (-15.8767, -17.998) arc (228.583:288.917:24) -- (13.0236, -35.3614)arc (118.011:187.497:10) arc (279.181:241.688:69.1) -- cycle;
\fill (10.5523, -21.5557) arc (296.083:311.417:24) -- (33.2593, -35.3806)arc (298.195:295.57:172.01) arc (31.6501:100.731:10) -- cycle;
\fill (17.998, -15.8767) arc (318.583:356.417:24) -- (71.3531, -1.5)arc (188.627:221.289:10) arc (312.84:298.995:193.11) -- cycle;
\fill[unicampred] (31.67, 75.56) circle (7);
\fill[unicampred] (81.24, 0) circle (7);
\fill[unicampred] (17.72, -44.19) circle (7);
\fill (-64, -80) -- (-64, -86.9041) .. controls (-63.9775, -89.3328) and (-63.9775, -89.3328) .. (-63.955, -89.8276) .. controls (-63.8426, -92.0315) and (-63.5952, -92.9535) .. (-62.8531, -93.7181) .. controls (-61.7286, -94.91) and (-59.9295, -95.2249) .. (-54.1949, -95.2249) .. controls (-50.7091, -95.2249) and (-49.2024, -95.1349) .. (-47.8531, -94.8651) .. controls (-46.1439, -94.5052) and (-44.997, -93.3133) .. (-44.7496, -91.6042) .. controls (-44.6147, -90.5472) and (-44.5922, -90.2999) .. (-44.5247, -86.9041) -- (-44.5247, -80) -- (-49.09, -80) -- (-49.09, -86.9041) .. controls (-49.09, -87.2189) and (-49.1124, -87.9385) .. (-49.1349, -88.4558) .. controls (-49.2024, -89.8726) and (-49.3148, -90.2774) .. (-49.6747, -90.6597) .. controls (-50.1694, -91.1994) and (-51.1139, -91.3343) .. (-54.2624, -91.3343) .. controls (-58.1529, -91.3343) and (-59.03, -91.087) .. (-59.2774, -89.8951) .. controls (-59.3898, -89.2879) and (-59.3898, -89.2654) .. (-59.4348, -86.9041) -- (-59.4348, -80) -- (-64, -80);
\fill (-41.2491, -80) -- (-41.2491, -95) -- (-36.5939, -95) -- (-36.7064, -83.8006) -- (-36.2116, -83.8006) -- (-28.3181, -95) -- (-20.3796, -95) -- (-20.3796, -80) -- (-25.0347, -80) -- (-24.9223, -91.1769) -- (-25.3945, -91.1769) -- (-33.2431, -80) -- (-41.2491, -80);
\fill (-17.0509, -80) -- (-17.0509, -95) -- (-12.2608, -95) -- (-12.2608, -80) -- (-17.0509, -80);
\fill (5.3068, -89.2204) .. controls (5.28431, -91.4468) and (5.21684, -91.4693) .. (0.201848, -91.4693) .. controls (-4.61075, -91.4693) and (-4.67821, -91.4243) .. (-4.67821, -87.5787) .. controls (-4.67821, -85.3073) and (-4.4983, -84.4303) .. (-3.95857, -83.9805) .. controls (-3.5088, -83.6207) and (-2.67671, -83.5307) .. (0.17936, -83.5307) .. controls (3.05792, -83.5307) and (4.38476, -83.6657) .. (4.72209, -83.958) .. controls (4.99195, -84.1829) and (5.03693, -84.3853) .. (5.08191, -85.3748) -- (9.64713, -85.3748) -- (9.64713, -84.6327) .. controls (9.64713, -82.2264) and (9.19735, -81.1244) .. (7.91549, -80.4723) .. controls (6.88101, -79.9325) and (5.48671, -79.7751) .. (1.48371, -79.7751) .. controls (-6.36487, -79.7751) and (-8.00655, -80.1574) .. (-8.81614, -82.1814) .. controls (-9.15348, -83.0135) and (-9.24343, -84.2504) .. (-9.24343, -87.6012) .. controls (-9.24343, -89.5127) and (-9.19845, -90.7496) .. (-9.13099, -91.4243) .. controls (-8.95108, -92.9535) and (-8.5013, -93.7856) .. (-7.51179, -94.3478) .. controls (-6.2974, -95.045) and (-4.40835, -95.2249) .. (1.46122, -95.2249) .. controls (4.78956, -95.2249) and (6.43124, -95.09) .. (7.64563, -94.7076) .. controls (9.30979, -94.1679) and (9.93948, -93.066) .. (9.93948, -90.6372) .. controls (9.93948, -90.4123) and (9.91699, -90.075) .. (9.8945, -89.2204) -- (5.3068, -89.2204);
\fill (29.0254, -95) -- (34.2203, -95) -- (26.3717, -80) -- (19.2653, -80) -- (11.3267, -95) -- (16.5441, -95) -- (17.826, -92.4588) -- (27.7435, -92.4588) -- (29.0254, -95) (26.1468, -89.1754) -- (19.4452, -89.1754) -- (22.3462, -83.4858) -- (23.2458, -83.4858) -- (26.1468, -89.1754);
\fill (36.1297, -80) -- (36.1297, -95) -- (40.6275, -95) -- (40.515, -83.7781) -- (41.3246, -83.7781) -- (46.9693, -95) -- (50.7699, -95) -- (56.4146, -83.7781) -- (57.1792, -83.7781) -- (57.0668, -95) -- (61.5645, -95) -- (61.5645, -80) -- (53.3336, -80) -- (48.8584, -89.7376) -- (44.4056, -80) -- (36.1297, -80);
\fill (64.896, -95) -- (69.4612, -95) -- (69.4612, -91.3343) -- (75.2183, -91.3343) .. controls (78.4567, -91.3343) and (79.4687, -91.2444) .. (80.4807, -90.907) .. controls (82.2798, -90.2549) and (82.887, -88.9505) .. (82.887, -85.6222) .. controls (82.887, -81.7766) and (82.0774, -80.5622) .. (79.2438, -80.1349) .. controls (78.4117, -80.0225) and (78.0069, -80) .. (75.1733, -80) -- (64.896, -80) -- (64.896, -95) (69.4612, -87.5787) -- (69.4612, -83.7556) -- (75.1733, -83.7556) .. controls (77.1748, -83.7781) and (77.2198, -83.7781) .. (77.5347, -83.8906) .. controls (78.1194, -84.093) and (78.3218, -84.5427) .. (78.3218, -85.6222) .. controls (78.3218, -86.6567) and (78.1868, -87.0615) .. (77.7596, -87.3313) .. controls (77.3772, -87.5562) and (77.2873, -87.5562) .. (75.1733, -87.5787) -- (69.4612, -87.5787);
\end{tikzpicture}}}%
\parbox[b]{\textwidth}{\centering\sffamily%
{\bfseries\fontsize{15.5pt}{1em}\selectfont\uppercase{Universidade Estadual de Campinas}}\\
\fontsize{11.3pt}{1.2em}\selectfont\uppercase{Faculdade de Engenharia Elétrica e de Computação}}

\vskip 5cm

{\large\scshape Nome do Autor\par}

\vskip 4cm

{\LARGE\bfseries Título da Dissertação ou Tese\\(na língua original)\par}

%\vskip 2cm
%
%{\LARGE\bfserie Título da Dissertação ou Tese\\(em português se necessário)\par}

\vfill

Campinas

20XX

\end{center}


\cleardoublepage
\begin{center}

{\large\scshape Nome do Autor\par}

\vfill

{\large\scshape Título da Dissertação ou Tese\\(na língua original)\par}

%\vskip 1cm
%
%{\large\sc Título da Dissertação ou Tese\\(em português se necessário)\par}

\vfill

\end{center}

\hfill\parbox{0.45\linewidth}{
% Para documento em português:
Dissertação/Tese apresentada à Faculdade de Engenharia Elétrica e de Computação da Universidade Estadual de Campinas como parte dos requisitos exigidos para a obtenção do título de Mestre(a)/Doutor(a) em <NOME DO TÍTULO>, na Àrea de <NOME DA ÁREA>.
%Em inglês ou espanhol caso o documento não seja redigido em português:
%Thesis/Dissertation presented to the Faculty/Institute of the University of Campinas in partial fulfillment of the requirements for the degree of Master/Doctor, in the area of <NAME OF AREA>.
% No caso de Cotutela Internacional de Tese, incluir a seguinte informação após o nome do Curso e da Área, se houver: "no âmbito do Acordo de Cotutela firmado entre a Unicamp e a <NOME DA UNIVERSIDADE (PAÍS)>"
}

\vfill

\begin{flushleft}

{\itshape Supervisor/Orientador:} <Nome do orientador>

\vskip 1.8ex

{\itshape Co-supervisor/coorientador:} <Nome do coorientador>

\vskip 4.2ex

\parbox{0.45\textwidth}{\small\scshape Este exemplar corresponde à versão final da dissertação/tese defendida pelo(a) aluno(a) <Nome do Aluno>, e orientada pelo(a) Prof(a). Dr(a). <Nome do Orientador>.}

\vskip 2cm

\end{flushleft}

\begin{center}

Campinas

20XX

\end{center}


\includepdf{ficha.pdf}

\cleardoublepage
\fronttitle{Comissão Julgadora -- Dissertação de Mestrado}
%\fronttitle{Comissão Julgadora -- Tese de Doutorado}

{\raggedright

\frontitem{Candidato:} Nome do aluno \frontitem{RA:} XXXXXX

\vskip 1.8ex

\frontitem{Data da defesa:} XX de xxxx de XXXX

\vskip 1.8ex

\frontitem{Título da Tese:} ``Título da dissertação ou tese.''

\vskip 1.8ex

Prof.\ Dr.\ Nome do Orientador (Presidente, {\small FEEC}/{\small UNICAMP})

Prof.\ Dr.\ Nome do Membro Externo (Instituição)

%Prof.\ Dr.\ Nome do Segundo Membro Externo (Instituição)

Prof.\ Dr.\ Nome do Membro Interno ({\small FEEC}/{\small UNICAMP})

%Prof.\ Dr.\ Nome do Segundo Membro Interno ({\small FEEC}/{\small UNICAMP})

\vskip 1.8ex

A ata de defesa, com as respectivas assinaturas dos membros da Comissão Julgadora, encontra-se no processo de vida acadêmica do aluno.

}


\cleardoublepage
\fronttitle{Dedicatória}


\cleardoublepage
\fronttitle{Agradecimentos}


\cleardoublepage
\fronttitle{Resumo}

Resumo redigido obrigatoriamente em português, contendo no máximo 500 palavras.

\cleardoublepage

\fronttitle{Abstract}

Resumo traduzido para o inglês.


\listoffigures

\listoftables

{
  \fancypagestyle{plain}{\fancyhf{}}
  \renewcommand{\headrulewidth}{0pt}
  \printnoidxglossary[
    type=acronym,
    title={Lista de Abreviaturas},
    nonumberlist=true,
    style=list,
  ]
  \thispagestyle{empty}
}

\chapter*{Lista de Símbolos}
\thispagestyle{empty}

\begin{description}[labelwidth=1cm]
\item[$c_0$] velocidade da luz no vácuo
\item[$x$] coordenada espacial
\end{description}



\tableofcontents

\clearpage

\pagestyle{plain}
\onehalfspacing

\glsresetall

\chapter{Introdução}

Texto introdutório.

O pacote \emph{siunitx} é utilizado para possibilitar a inserção de números e unidades corretamento formatados: \SI{-1.2e-3}{\celsius/m^2}.
A formatação também funciona em equações, como mostrado em~(\ref{eq:angles}).

\begin{equation}
\label{eq:angles}
\SI{2\pi}{rad} = \SI{360}{\degree}
\end{equation}

Aproveitamos para fazer refência à fig.~\ref{fig:logo} e às fontes~\cite{oetiker_not_2015, latex_wikibook}.

\begin{figure}[htpb]
\centering
\includegraphics{unicamp-logo}
\caption{Logotipo da {\small UNICAMP}}
\label{fig:logo}
\end{figure}

Por fim um exemplo de tabela limpa, sem poluição visual, pode ser visto na tabela~\ref{t:xxx}.

\begin{table}[hbpt]
\centering
\caption{Legenda da tabela}
\label{t:xxx}
\begin{tabular}{lr}
\toprule
Constante & Valor \\
\midrule
$\pi$ & \num{3.14159265}\ldots \\
$c_0$ &	\SI{299792458}{m/s} \\
\bottomrule
\end{tabular}
\end{table}

\section{Organização da tese}

Primeiro a introdução e por fim a conclusão.



\chapter{Conclusão}



\printbibliography

\appendix

\chapter{Notas complementares}



\end{document}

